\begin{trapbox}
	\begin{description}[style=nextline, leftmargin=2.8em, labelsep=0.5em, itemsep=0.35em]
		\item[\textbf{\textcolor{CTrap}{易错点一（混淆对立与互斥）}}] 错误认为若两事件互斥，则它们互为对立事件。
		\item[\textbf{\textcolor{CTrap}{正确理解（概念区分）：}}] 对立事件必须满足互斥且并集为样本空间，而互斥事件只要求互斥，不要求并集为样本空间。
		\item[\textbf{\textcolor{CTrap}{错因分析（概念掌握不足）：}}] 对事件关系的基本概念理解不准确，将互斥当作对立。
	\end{description}

	\medskip
	\begin{description}[style=nextline, leftmargin=2.8em, labelsep=0.5em, itemsep=0.35em]
		\item[\textbf{\textcolor{CTrap}{易错点二（忽略对立条件）}}] 未确认事件是否真正对立就直接使用公式 $P(A)=1-P(B)$，比如 $A$ 与 $B$ 仅是互斥但不完备。
		\item[\textbf{\textcolor{CTrap}{正确理解（验证完备性）：}}] 使用公式前必须确认 $A$ 和 $B$ 满足 $A\cup B=\Omega$ 且 $A\cap B=\emptyset$。
		\item[\textbf{\textcolor{CTrap}{错因分析（盲目套公式）：}}] 忽略公式的前提条件，随意套用导致错误。
	\end{description}

	\medskip
	\begin{description}[style=nextline, leftmargin=2.8em, labelsep=0.5em, itemsep=0.35em]
		\item[\textbf{\textcolor{CTrap}{易错点三（计算反面概率出错）}}] 在计算 $P(\overline{A})$ 时遗漏样本点，如古典概型中计数错误。
		\item[\textbf{\textcolor{CTrap}{正确理解（仔细计数）：}}] 列出 $\overline{A}$ 包含的所有可能结果，确保无重复、无遗漏。
		\item[\textbf{\textcolor{CTrap}{错因分析（计数不完整）：}}] 样本空间较大时手工枚举易漏，应使用组合计数或系统列举。
	\end{description}
\end{trapbox}
